\section*{الفصل \ref{ch:b3:gravity-geodesy} --- الجاذبية وشكل الأرض}

\begin{solution}{exo:b3:gravity-geodesy:1}
$b = a(1 - f) = 6378137\times0.99664719 = \qty{6356752.3}{m}$،
و$a - b = \qty{21384.7}{m}$، أي نحو \qty{21.4}{km}. كليرو:
$f = (3J_2 + m)/2 = (0.00324789 + 0.00344979)/2 = \num{0.0033488}$، أو
$1/298.61$، مقابل $1/298.26$: أصغر بمقدار \qty{0.12}{\percent}، أي بقدر
حدود الرتبة الثانية التي تُهمِلها العلاقة.
\end{solution}

\begin{solution}{exo:b3:gravity-geodesy:2}
$\sigma = \qty{6.0e13}{kg}/\qty{5.0e11}{m^2} = \qty{120}{kg.m^{-2}}$، إذن
$h_w = 120/1000 = \qty{0.12}{m}$: \qty{12}{cm} من الماء. وعند تصريفه في
المحيط، $60/360 = \qty{0.17}{mm}$ من مستوى سطح البحر.
\end{solution}

\begin{solution}{exo:b3:gravity-geodesy:3}
عند خط الاستواء $\gamma = \qty{9.780327}{m.s^{-2}}$. وعند \ang{45}،
$\sin^2\varphi = 0.5$ و$\sin^22\varphi = 1$: $\gamma = 9.780327\times1.0026454
= \qty{9.806200}{m.s^{-2}}$. وعند القطبين، $9.780327\times1.0053024 =
\qty{9.832186}{m.s^{-2}}$. الفرق هو \qty{0.051859}{m.s^{-2}}، أو
\qty{5186}{\milli\gal}.
\end{solution}

\begin{solution}{exo:b3:gravity-geodesy:4}
$D = ET_e^3/(12(1 - \nu^2)) = 10^{11}T_e^3/11.25$: \qty{8.89e21}{N.m} من أجل
\qty{10}{km}، \qty{2.40e23}{N.m} من أجل \qty{30}{km}. مع $(\rho_m - \rho_w)g =
\qty{22269}{N.m^{-3}}$، $\alpha = (4D/22269)^{1/4}$: \qty{35.5}{km}
و\qty{81.0}{km}. تنتهي الحفرة المحيطة عند $3\pi\alpha/4$، \qty{84}{km} و\qty{191}{km}؛
تبلغ قمم \omterm{def:b3:gravity-geodesy:alpha}{الانتفاخ الأمامي} عند $\pi\alpha$، \qty{112}{km} و\qty{255}{km}.
\end{solution}

\begin{solution}{exo:b3:gravity-geodesy:5}
(أ) $T = -\gamma N = -9.80\times50 = \qty{-490}{J.kg^{-1}}$. (ب) يميل الجيود
بمقدار $10/100\,000 = \num{1.0e-4}$ راديان بالنسبة إلى الإهليلجي، ويميل الشاقول
وهو العمودي على الجيود، بالزاوية نفسها بالنسبة إلى عمودي الإهليلجي:
$\num{1.0e-4}\times206\,265 = 20.6''$.
\end{solution}

\begin{solution}{exo:b3:gravity-geodesy:6}
$H = J_2/(C/Ma^2) = \num{1.08263e-3}/0.3308 = \num{0.0032728}$،
و$1/H = 305.55$. $J_2$ هو نفسه بالفرض؛ ولو كانت الأرض متجانسة لكان
$C/(Ma^2) = 0.4$ ومنه $H = J_2/0.4 = \num{0.0027066}$: لكانت المبادرة
أبطأ. إنّ قياس $J_2$ و$H$ معًا هو ما يكشف تركّز الكتلة
نحو المركز.
\end{solution}

\begin{solution}{exo:b3:gravity-geodesy:7}
$r = \qty{6878.137}{km}$؛ $n = \sqrt{GM/r^3} = \qty{1.10678e-3}{rad.s^{-1}}$، أي
دور قدره $2\pi/n = \qty{94.6}{min}$. $\dot\Omega = -1.5\times\num{1.10678e-3}
\times\num{1.08263e-3}\times(6378.137/6878.137)^2\times\cos\ang{89} =
\qty{-2.70e-8}{rad.s^{-1}}$، أو \ang{-0.134} في اليوم، غربًا. يمرّ مدار شبه قطبي
فوق كل خط عرض، فتُرسم الأرض كلها؛ كما يتقدّم المدار
ببطء، فتملأ مسارات السطح الفراغات بانتظام.
\end{solution}

\begin{solution}{exo:b3:gravity-geodesy:8}
$e^{-2\pi\times500/100} = e^{-31.4} = \num{2.3e-14}$: \qty{2e-13}{\milli\gal}، أي لا شيء.
$e^{-2\pi\times500/1000} = e^{-3.14} = 0.0432$:
\qty{0.43}{\milli\gal}.
\end{solution}

\begin{solution}{exo:b3:gravity-geodesy:9}
$3\pi\alpha/4 = \qty{120}{km}$ يعطي $\alpha = \qty{50.9}{km}$.
$D = (\rho_m - \rho_w)g\alpha^4/4 = 22269\times(50\,930)^4/4 =
\qty{3.75e22}{N.m}$، و$T_e = (11.25D/E)^{1/3} = \qty{16.2}{km}$. تعبر الصفيحة
مستواها غير المنثني عند بُعد $\pi\alpha/2 = \qty{80}{km}$ من الخندق.
\end{solution}

\begin{solution}{exo:b3:gravity-geodesy:10}
(أ) بوضع $u = x/\alpha$، $\mathrm d^4/\mathrm dx^4 = \alpha^{-4}\,\mathrm
d^4/\mathrm du^4$. بكتابة $\cos u + \sin u = \sqrt2\cos(u - \pi/4)$، فتكون الدالة
هي $\sqrt2w_0\,\mathrm{Re}\,e^{(-1 + i)u - i\pi/4}$،
و$(-1 + i)^4 = -4$: إذن $w'''' = -4w/\alpha^4$. ثم $Dw'''' + (\rho_m -
\rho_i)gw = (-4D/\alpha^4 + (\rho_m - \rho_i)g)w = 0$، لأنّ $\alpha^4 =
4D/((\rho_m - \rho_i)g)$. (ب) $w'' = -(2w_0/\alpha^2)e^{-u}(\cos u - \sin u)$، إذن
$M = -Dw'' = (2Dw_0/\alpha^2)e^{-u}(\cos u - \sin u)$: أكبر ما يكون في المقدار المطلق عند
$u = 0$، وينعدم حيث $\tan u = 1$، أي $x = \pi\alpha/4$. (ج) تحت الحمل
تكون الصفيحة مقعّرة نحو الأعلى، وسطحها العلوي في انضغاط؛ وبعد $\pi\alpha/4$
تصبح محدّبة نحو الأعلى، وسطحها العلوي في تمدّد. وهناك $M' \propto
-e^{-u}\cos u$ ينعدم عند $u = \pi/2$: يبلغ التمدّد أقصاه عند
$x = \pi\alpha/2$، على المنحدر الخارجي للحفرة المحيطة نحو \omterm{def:b3:gravity-geodesy:alpha}{الانتفاخ الأمامي}، حيث
$|M|$ هو $e^{-\pi/2} = \qty{21}{\percent}$ من قيمته تحت الحمل. وهذا هو
الموضع الذي تفتح فيه الصدوع العادية والشقوق في أعلى الصفيحة.
\end{solution}

\begin{solution}{exo:b3:gravity-geodesy:11}
$\sqrt{1 + \eta} = 2.5(1 - 1.5\times0.3308) = 1.2595$، إذن $1 + \eta = 1.5863$
و$\tfrac52m/f = 2.5863$: $f_h = 2.5\times0.00344979/2.5863 = \num{0.0033346}$،
$1/f_h = 299.9$. والفرق $f - f_h = \num{1.82e-5}$؛ وبضربه في $a$، يرتفع
خط الاستواء بمقدار \qty{116}{m} فوق الشكل الهيدروستاتيكي. والأرض ليست تمامًا
مائعًا في حالة توازن: إذ يتحمّل باطنها إجهادات، ناتجة عن متانة
الغلاف الصخري أو عن فروق التدفّق والكثافة في الوشاح، تحفظ
الشكل مبتعدًا عنه قليلًا.
\end{solution}

\begin{solution}{exo:b3:gravity-geodesy:12}
يجب أن تتقدّم العقدة بمقدار $2\pi/(365.242\times86\,400) =
\qty{1.991e-7}{rad.s^{-1}}$. عند $r = \qty{7078.137}{km}$، $n =
\qty{1.0602e-3}{rad.s^{-1}}$ و$\tfrac32nJ_2(a/r)^2 = \num{1.398e-6}$، إذن
$\cos i = -\num{1.991e-7}/\num{1.398e-6} = -0.1424$ و$i = \ang{98.2}$: مدار
قليل الارتداد، متزامن مع الشمس.
\end{solution}

\begin{solution}{pb:b3:gravity-geodesy:1}
\textbf{1.} $m = (\num{7.292115e-5})^2\times6378137^2\times6356752.3/
\num{3.986005e14} = \num{0.00344979}$.

\textbf{2.} $f = (3J_2 + m)/2 = \num{0.0033488}$, $1/f = 298.61$.

\textbf{3.} $\sqrt{1 + \eta} = 2.5(1 - 1.5\times0.3308) = 1.2595$,
$\tfrac52m/f_h = 2 + 0.5863$، ومنه $f_h = \num{0.0033346}$ و$1/f_h = 299.88$.

\textbf{4.} قيمة كليرو، $1/298.61$، تقع ضمن \qty{0.12}{\percent} من
القيمة المرصودة $1/298.26$؛ أمّا القيمة الهيدروستاتيكية، $1/299.88$، فهي أصغر بمقدار
$\Delta f = \num{1.82e-5}$، و$a\,\Delta f = \qty{116}{m}$: يعلو خطّ الاستواء
\qty{116}{m} فوق الشكل الهيدروستاتيكي.

\textbf{5.} $H = \num{1.08263e-3}/0.3308 = \num{0.0032728}$ ($1/305.6$).

\textbf{6.} $\tfrac12\times\qty{80}{km}\times\qty{6}{km} = \qty{240}{km^2} =
\qty{2.40e8}{m^2}$.

\textbf{7.} $V_0 = (2800 - 1030)\times9.81\times\num{2.40e8} =
\qty{4.17e12}{N.m^{-1}}$.

\textbf{8.} إيري: $(\rho_m - \rho_v)r = (\rho_v - \rho_w)h$، ومنه $r =
6\times1770/500 = \qty{21.2}{km}$ تحت القمّة.

\textbf{9.} ينحني قاع البحر انحناءً كبيرًا إلى ما بعد قاعدة الظهرة:
الحمل لا يُعوَّض عمودًا فعمودًا، بل بانثناء الصفيحة عبر
منطقة أوسع أربع مرات. للصفيحة صلابة: \omterm{def:b3:gravity-geodesy:regional}{التعويض الأيزوستاتي الإقليمي}.

\textbf{10.} $\alpha = 165/(3\pi/4) = \qty{70.0}{km}$.

\textbf{11.} $\alpha = 220/\pi = \qty{70.0}{km}$: القيمة نفسها، كما تقتضيه صفيحة
مرنة واحدة.

\textbf{12.} $D = (\rho_m - \rho_w)g\alpha^4/4 = 2270\times9.81\times(\num{7.00e4})^4/4
= \qty{1.34e23}{N.m}$.

\textbf{13.} $T_e = (12(1 - \nu^2)D/E)^{1/3} = (11.25\times\num{1.34e23}/
\num{1e11})^{1/3} = \qty{24.7}{km}$.

\textbf{14.} $w_0 = V_0\alpha^3/(8D) = \num{4.17e12}\times(\num{7.00e4})^3/
(8\times\num{1.34e23}) = \qty{1334}{m}$.

\textbf{15.} $w_0e^{-\pi} = 1334\times0.0432 = \qty{58}{m}$.

\textbf{16.} $2\times1334\times\num{7.00e4}\times2270\times9.81 =
\qty{4.16e12}{N.m^{-1}}$، يساوي $V_0$ ضمن حدود تقريب $D$.

\textbf{17.} \qty{24.7}{km} تقع ضمن نطاق ظهرة هاواي. الصفيحة
تغوص \qty{1.3}{km} تحت المحور بدلًا من أن ينمو لها جذر بسمك
\qty{21}{km}: إذ يتوزّع الطفو نفسه على \qty{330}{km} من الخندق المحيط.

\textbf{18.} $\lambda = 2\pi\times\qty{70.0}{km} = \qty{440}{km}$,
$k = 1/\alpha = \qty{1.43e-5}{m^{-1}}$.

\textbf{19.} $\sigma = (2900 - 1030)\times500 = \qty{9.35e5}{kg.m^{-2}}$، ومنه
$2\pi G\sigma = \qty{3.92e-4}{m.s^{-2}} = \qty{39.2}{\milli\gal}$ عند قاع البحر؛
مضروبًا في $e^{-k\times\qty{5}{km}} = 0.931$، \qty{36.5}{\milli\gal} عند مستوى سطح البحر
(سالبة حيث تُدفع الصفيحة إلى أسفل).

\textbf{20.} $\sigma = (3300 - 2900)\times500 = \qty{2.0e5}{kg.m^{-2}}$:
\qty{8.39}{\milli\gal}، مضروبًا في $e^{-k\times\qty{12}{km}} = 0.843$،
\qty{7.07}{\milli\gal}. ويكون المجموع \qty{43.6}{\milli\gal}.

\textbf{21.} $e^{-k\times\qty{400}{km}} = e^{-5.71} = 0.0033$: \qty{0.14}{\milli\gal}.

\textbf{22.} $e^{-2\pi\times400/80} = e^{-31.4} = \num{2.3e-14}$: القمر الصناعي
لا يرى من الظهرة ذاتها شيئًا، بل الانثناء الواسع فحسب، وحتى هذا
يضعف بمقدار ثلاثمئة مرّة؛ ويجب قياس الشذوذ الخاص بالظهرة في البحر.

\textbf{23.} $n \approx 2\pi\times6371/440 = 91$.

\textbf{24.} من المسائل 10 إلى 13، \omterm{def:b3:gravity-geodesy:flexure}{السمك المرن الفعّال} للصفيحة
تحت ظهرة أورو هو
\textbf{\qty[text-series-to-math=true]{24.7}{km}}.
\end{solution}
